\section{Bližší vymezení strojového učení}\label{sec:ai-vymezeni-ml}

U~běžného algoritmu stanoví programátor přesný postup výpočtu. U~strojového učení místo toho určíme třídu modelů, způsob měření chyby a~postup, kterým se parametry modelu upravují podle dat. Výsledné konkrétní pravidlo tedy nevzniká pouze přímým zápisem programu, ale také trénováním.

Učení lze popsat pomocí tří složek:
\begin{itemize}
    \item \textbf{úloha} říká, co má model provádět, například předpovídat cenu nebo rozpoznávat spam;
    \item \textbf{zkušenost} představuje data nebo interakce, z~nichž model čerpá;
    \item \textbf{měřítko úspěchu} určuje, podle čeho poznáme zlepšení modelu.
\end{itemize}
O~učení hovoříme tehdy, když se s~rostoucí zkušeností zlepšuje výkon modelu v~dané úloze podle zvoleného měřítka.

\subsection{Trénování a~použití modelu}

\emph{Trénováním} rozumíme hledání parametrů modelu na základě dat. Při následném \emph{použití} jsou parametry pevné a~model vypočítává výstupy pro nové vstupy. Tyto dvě fáze je nutné odlišovat: náročné hledání parametrů může probíhat jednou, zatímco hotový model potom používáme mnohokrát.

\begin{figure}[ht]
    \centering
    \begin{tikzpicture}[every node/.style={font=\normalsize}]
    \node[task,fill=blue!10,minimum width=28mm] (train) at (0,1.2)
        {trénovací data};
    \node[task,fill=orange!15,minimum width=28mm] (learn) at (4.8,1.2)
        {trénování};
    \node[task,fill=green!12,minimum width=28mm] (model) at (9.6,1.2)
        {hotový model};
    \node[task,fill=blue!10,minimum width=28mm] (new) at (4.8,-1.1)
        {nový vstup};
    \node[task,fill=green!12,minimum width=28mm] (pred) at (9.6,-1.1)
        {předpověď};
    \draw[diredge] (train) -- (learn);
    \draw[diredge] (learn) -- node[above]{parametry} (model);
    \draw[diredge] (new) -- (pred);
    \draw[diredge] (model) -- (pred);
\end{tikzpicture}

    \caption{Odlišení trénování modelu od jeho následného použití.}
    \label{fig:ai-trenovani-pouziti}
\end{figure}

Pro kontrolu schopnosti zobecnit rozdělujeme dostupná data na tři části:
\begin{itemize}
    \item \textbf{trénovací data} slouží k~určení parametrů modelu;
    \item \textbf{validační data} slouží k~volbě nastavení, například hodnoty \(k\), rozhodovacího prahu nebo počtu epoch;
    \item \textbf{testovací data} použijeme až na konci pro nezávislé závěrečné vyhodnocení.
\end{itemize}
Testovací data nesmějí ovlivnit trénování ani ladění. Opakované rozhodování podle výsledku na testovací množině by z~ní postupně udělalo další validační množinu.

\subsection{Základní členění}

\begin{figure}[ht]
    \centering
    \begin{tikzpicture}[every node/.style={font=\normalsize}]
    \node[task,fill=orange!15,minimum width=33mm] (ml) at (5,1.6)
        {strojové učení};
    \node[task,fill=blue!10,minimum width=37mm] (sup) at (0,-0.6)
        {učení s~učitelem\\známé výstupy};
    \node[task,fill=green!12,minimum width=37mm] (unsup) at (5,-0.6)
        {učení bez učitele\\pouze vstupy};
    \node[task,fill=red!10,minimum width=37mm] (rl) at (10,-0.6)
        {zpětnovazební učení\\odměny};
    \draw[diredge] (ml) -- (sup);
    \draw[diredge] (ml) -- (unsup);
    \draw[diredge] (ml) -- (rl);
\end{tikzpicture}

    \caption{Tři základní druhy strojového učení.}
    \label{fig:ai-druhy-uceni}
\end{figure}

\begin{itemize}
    \item Při \textbf{učení s~učitelem} známe u~trénovacích objektů požadované výstupy.
    \item Při \textbf{učení bez učitele} výstupy nebo třídy předem dány nejsou a~hledáme strukturu ve vstupních datech.
    \item Při \textbf{zpětnovazebním učení} agent volí akce, pozoruje jejich následky a~učí se z~odměn.
\end{itemize}

\warningbox{Nízká chyba na trénovacích datech sama o~sobě nestačí. Smyslem modelu je pracovat s~novými daty. Pokud si model pouze zapamatuje zvláštnosti trénovací množiny, dochází k~přeučení.}
